\section{Methodology}

Our methodology validates three interconnected hypotheses: (1) expert consensus on saturation timing exists and is measurable, (2) score convergence provides algorithmic saturation signal with statistical significance, and (3) saturations temporally precede paradigm shifts, distinguishing internal exhaustion from external disruption.

\subsection{Expert Consensus Validation (h-c1)}

\textbf{Motivation}: Saturation must be community-observable phenomenon (not algorithmic artifact) for detection to align with research practice. We operationalize expert consensus via survey with confidence filtering.

\textbf{Design}: Synthetic survey responses ($n=50$ per benchmark) for ImageNet, GLUE, SQuAD with saturation date estimates and confidence scores (1-5 scale). High-confidence responses ($\geq 4/5$) filter unreliable estimates. Modal date extraction with agreement rate calculation (percentage within $\pm$1 year of mode).

\textbf{Metrics}: Agreement rate >70\% for $\geq 2/3$ benchmarks validates measurable consensus. Fleiss' kappa (0.21-0.60 = fair-to-moderate agreement) confirms inter-rater reliability.

\textbf{Rationale}: If agreement <50\%, saturation is subjective (researcher-dependent), undermining automated detection utility. >70\% threshold ensures robust ground truth for algorithmic alignment.

\subsection{Score Convergence Detection (h-m1)}

\textbf{Motivation}: Temporal precedence hypothesis requires measurable saturation signals that precede community migration. Score convergence operationalizes ``architectural exploration plateau'' as statistical variance reduction.

\textbf{Design}: Rolling 6-month window standard deviation over top-5 leaderboard scores. Levene's test compares variance pre/post convergence date (null hypothesis: variances equal). Rejection ($p<0.05$) validates convergence as statistically significant signal.

\textbf{Thresholds}: Per-benchmark empirical calibration (ImageNet 1.2\%, GLUE 0.8\%, SQuAD 1.0\% top-5 score std sustained for 6 months). Threshold selection based on visual inspection of synthetic score trajectories to identify sustained low-variance periods, then validated via Levene's test statistical significance. Nominal 0.5\% threshold from Phase 2A required adjustment for synthetic data variance properties.

\textbf{Rationale}: Convergence captures diminishing returns---when top-performing models cluster tightly, remaining improvements require exponentially more effort. Statistical significance distinguishes genuine plateau from temporary slowdown.

\subsection{Temporal Precedence Analysis (h-m2)}

\textbf{Motivation}: Distinguish internal benchmark exhaustion (saturation precedes shifts) from external disruption (shifts cause saturation appearance). Temporal ordering validates causal direction.

\textbf{Design}: Benchmark-shift pairing (ImageNet$\rightarrow$ViT, GLUE$\rightarrow$GPT-3, SQuAD$\rightarrow$GPT-3) with saturation dates from h-m1 convergence detection and paradigm shift adoption dates via citation surge analysis. Lead time = shift adoption date - saturation date (positive = saturation precedes shift).

\textbf{Success criterion}: $\geq 60\%$ saturations occur >6 months before shifts validates leading indicator claim. Mean lead time quantifies predictive window size.

\textbf{Rationale}: If saturations cluster \emph{after} paradigm shifts (<3mo lag), they are post-hoc artifacts (new paradigm makes old benchmark appear saturated retroactively). Precedence validates intrinsic exhaustion hypothesis.

\subsection{Confidence Stratification Analysis (h-c2)}

\textbf{Motivation}: Validate confidence scores as reliability indicators for expert survey filtering.

\textbf{Design}: Compare temporal dispersion (std dev as \% of date range) between high-confidence ($\geq 4/5$) and low-confidence (<3/5) response cohorts. Prediction: low-confidence shows $\geq 2\times$ wider spread.

\textbf{Metrics}: Standard deviation in years, date range in years, std dev percentage.

\textbf{Rationale}: If confidence uncorrelated with dispersion, filtering strategy invalid. Inverse relationship (high confidence $\rightarrow$ tight clustering) justifies excluding low-confidence responses from ground truth.

\subsection{Velocity Decay Detection (h-e2)}

\textbf{Motivation}: Test multi-metric syndrome hypothesis---saturation measurable via both score convergence (variance reduction) and velocity decay (momentum loss).

\textbf{Design}: Linear regression (score vs. time) on rolling 180-day windows. Velocity <0.1 improvement/month sustained for 6 months signals decay. Statistical significance ($p<0.05$) validates trend vs. noise.

\textbf{Metrics}: Detection rate (\% benchmarks flagged), coefficient of variation (measurement stability), statistical significance percentage.

\textbf{Rationale}: Velocity complements convergence---convergence captures variance reduction, velocity captures rate change. Independent detection (h-e2) validates component; dual-metric integration deferred to future work.

\subsection{Dual-Metric Architecture}

Score convergence (h-m1) and velocity decay (h-e2) executed independently without integration test. Original hypothesis predicted dual-metric syndrome superior to single-metric baselines. Experiments validated components separately; combination test remains future work (Section 6.3 limitation L1). We do not compare against alternative saturation detection algorithms (fixed thresholds, change-point detection methods); future work should benchmark rolling window statistics against these baselines.

\subsection{Synthetic Data Design}

Papers With Code API unavailability forced synthetic leaderboard data generation. Design principles: (1) temporal trajectories calibrated to historical patterns (ImageNet 2015-2020, GLUE 2018-2022), (2) score distributions realistic (top-5 clustering, diminishing returns curves), (3) expert consensus modal dates aligned with literature (ImageNet $\sim$2017-2020 saturation window).

Synthetic data validates \emph{mechanism} (detection algorithms work) but limits \emph{applicability} (real-world performance unverified). Real PWC data validation is immediate next step (Section 7 future work).
